\documentclass[10pt,a4paper]{amsart}
\usepackage[margin=19mm]{geometry}
\usepackage{amsmath,amssymb,mathtools}
\usepackage[scr=rsfs,cal=euler]{mathalfa}
\usepackage{xfrac}
\usepackage[unicode,hidelinks,bookmarks=false]{hyperref}
\newcommand{\ie}{i.e.\ }
\newcommand{\cf}{cf.\ }
\newcommand{\resp}{respectively\ }
\newcommand{\Subsection}{\subsection}
\providecommand{\nobreakdash}{\nobreak\hbox{-}}
\setlength{\emergencystretch}{2em}
\multlinegap0pt
\usepackage{bm}
\usepackage{bbm}
\usepackage{dsfont}
\usepackage{xspace}
\usepackage{cancel}
\usepackage{mathtools}
\usepackage{cleveref}
\usepackage[shortlabels]{enumitem}
\usepackage{amsthm}
\makeatletter
\@ifundefined{definition}{\newtheorem{definition}{Definition}}{}
\@ifundefined{theorem}{\newtheorem{theorem}{Theorem}}{}
\@ifundefined{claim}{\newtheorem{claim}[theorem]{Claim}}{}
\@ifundefined{lemma}{\newtheorem{lemma}[theorem]{Lemma}}{}
\@ifundefined{observation}{\newtheorem{observation}[theorem]{Observation}}{}
\@ifundefined{proposition}{\newtheorem{proposition}[theorem]{Proposition}}{}
\makeatother
\newcommand{\Pow}{{q_\star}} %Notaiton for (large) powers of rooted graphs
\newcommand{\mon}{\mathrm{mon}}
\newcommand{\pow}{q} %Notation for (small) powers of rooted graphs
\newcommand{\sm}{\setminus}
\newcommand{\sub}{\subseteq}
\renewcommand{\c}[1]{\mathcal{#1}}
\title[Rational Exponents for General Graphs]{Rational Exponents for General Graphs}
\author{Sean English, Sam Spiro}
\date{Source: arXiv:2506.19061v1 (CC BY 4.0)}
\begin{document}
\maketitle

\noindent\emph{Source and licence.} Rational Exponents for General Graphs. Version arXiv:2506.19061v1 (CC BY 4.0), distributed under the Creative Commons Attribution 4.0 International licence, as recorded in the arXiv licence metadata for this exact version and shown on the abstract page https://arxiv.org/abs/2506.19061v1. Full rights notice: licenses/arxiv-2506.19061-CC-BY-4.0.txt.\par\smallskip

\noindent\textit{Editorial scope.} Counted unit: the Proposition whose source label is \texttt{proposition:HellyTechnical} in the article (source0.tex lines 1175--1177, source label \texttt{proposition:HellyTechnical}), delivered together with the article's own complete original proof of it (source0.tex lines 1180--1228, one \texttt{proof} environment of 49 lines). No mathematics is written, repaired or re-derived: statement and proof are the article's own text line for line. The proof is detached from the statement in the source: the article states the result at lines 1175--1177 and prints its proof at lines 1180--1228, and the proof environment's own header names the statement, so the association is the article's own. Required context retained uncounted: observation monomorphisms same as copies (lines 226--230); lemma:canonicalReduction (lines 1013--1015); lemma:sunflower (lines 1022--1024); theorem:treeBoundEG (lines 1099--1101); thm:Helly (lines 1134--1136); condition highly extendable definition (a) leaf of T' (lines 1163--1171); condition highly extendable definition (c) many canonical morphisms (lines 1165--1169); condition highly extendable defintion (b) edges incident (lines 1165--1169). Every definition, display and label of those lines is retained unchanged. Omitted from this package and not redistributed: 1--225, 231--1012, 1016--1021, 1025--1098, 1102--1133, 1137--1162, 1170--1174, 1178--1179, 1229--1349. Nothing inside a retained excerpt is omitted or altered: every retained body is delivered exactly as the article's lines are written, so no omission receipt of any kind accompanies this package. Citations: the retained excerpts carry no citation command at all, so no bibliography entry of the article is transcribed and no reference is redistributed. Reference adaptation: none was needed and none was made; every reference inside the delivered unit resolves inside it, so no unresolved reference or citation is produced. Printed slips kept literal: no printed word, symbol, hypothesis or number of the counted statement or of its proof was altered. Layout and class: the article's own class is \texttt{article} with packages including amssymb, amsfonts, amscd, amsthm, xy, bbm, mathrsfs, tikz, geometry, mathtools, hyperref, cleveref, color, graphicx; the wrapper is the lane's pinned \texttt{amsart} editorial wrapper at 10pt on A4, which restates the theorem-like environments the retained text needs and adds the article's own macro definitions that the retained text needs (\texttt{\textbackslash Pow}, \texttt{\textbackslash c}, \texttt{\textbackslash mon}, \texttt{\textbackslash pow}, \texttt{\textbackslash sm}, \texttt{\textbackslash sub}), with extra wrapper packages (bm, bbm, dsfont, xspace, cancel, mathtools, cleveref, enumitem). The delivered numbering is the unit's own. Assets: no retained span contains a figure environment, a graphics-inclusion command, a PDF-page inclusion, a table or a TikZ picture; no image file is copied into this package, the package declares no asset dependency and no image file is redistributed with it. Licence: version arXiv:2506.19061v1 of this article is distributed under the Creative Commons Attribution 4.0 International licence, as recorded in the arXiv licence metadata for this exact version and shown on the abstract page https://arxiv.org/abs/2506.19061v1. A full rights notice is delivered at licenses/arxiv-2506.19061-CC-BY-4.0.txt. Characters: every retained excerpt is pure ASCII, so the delivered engine, XeLaTeX with its default Latin Modern text font, reports no missing character.
	\begin{observation}\label{observation monomorphisms same as copies}
		For every pair of graphs $H,G$, the number of copies of $H$ in $G$ equals
		\[\frac{\mon(H,G)}{\mathrm{aut}(H)},\]
		where $\mathrm{aut}(H)$ denotes the size of the automorphism group of $H$, i.e.\ the number of isomorphisms from $H$ to itself.
	\end{observation}

	\begin{lemma}\label{lemma:canonicalReduction}
		If $G$ and $H$ are graphs and if $\Phi\subseteq \mathrm{Mon}(H,G)$ is a collection of monomorphisms, then there exists a partition of $V(G)$ into sets $\{V_{\hat{x}}:\hat{x}\in V(H)\}$ and a subcollection $\Phi'\sub \Phi$ with $|\Phi'|\ge v(H)^{-v(H)}|\Phi|$ such that $\phi(\hat{x})\in V_{\hat{x}}$ for all $\hat{x}\in V(H)$ and $\phi\in \Phi'$.
	\end{lemma}

	\begin{lemma}\label{lemma:sunflower}
		For every graph $H$ and integer $\pow \ge 2$, there exists some integer $\Pow=\Pow(H,q)$ such that $\Pow\ge \pow$ and such that if $G$ is a graph and if $\Phi\subseteq \mathrm{Mon}(H,G)$ satisfies $|\Phi|\ge \Pow$, then there exists a set $R\subsetneq V(H)$ and distinct maps $\phi_1,\ldots,\phi_\pow\in \Phi$ such that for all $i\ne j$, we have $\phi_i(\hat{x})=\phi_j(\hat{y})$ if and only if $\hat{x}=\hat{y}$ and $\hat{x}\in R$.
	\end{lemma}

	\begin{theorem}\label{theorem:treeBoundEG}
		If $T\ne K_1$ is a tree and if $k,\pow\ge 2$ are integers, then any $n$-vertex graph $G$ which is $\c{F}_{T,k}^\pow$-free contains at most $O(e(G)^{k-1})$ copies of $T$. 
	\end{theorem}

	\begin{theorem}\label{thm:Helly}
		Let $k\geq 1$ be an integer. If $T$ is a tree and if $\c{T}$ is a $k$-Helly collection of leaf-cuttable subtrees, then there exist a set $E$ of at most $k-1$ edges of $T$ such that every $T_i\in \c{T}$ contains at least one edge of $E$.
	\end{theorem}

	\begin{definition}
		Let $T$ be a tree, $G$ a graph, $\{V_{\hat{x}}:\hat{x}\in V(T)\}$ a partition of $V(G)$ indexed by $V(T)$, and $k,\pow\ge 2$ a pair of integers.  We say that a subtree $T'\sub T$ is a \textbf{highly-extendable} with respect to some set $S\sub V(G)$ if the following conditions hold:
		\begin{enumerate}[label=(\alph*)]
			\item Every vertex of $\widehat{S}\sub V(T)$ is a leaf of $T'$,\label{condition highly extendable definition (a) leaf of T'}
			\item Every edge of $E(T)\sm E(T')$ which is incident to a vertex of $T'$ is incident to a vertex of $\widehat{S}$, \label{condition highly extendable defintion (b) edges incident}
			\item We have $|\Psi(T';S)|\geq \Pow(T',q)$ where $\Pow(T',q)$ is the real number guaranteed by \Cref{lemma:sunflower}.\label{condition highly extendable definition (c) many canonical morphisms}
		\end{enumerate}
		For each canonical monomorphism $\phi:V(T)\to V(G)$, we define $\c{T}_\phi$ to be the set of all subtrees $T'$ which are highly-extendable with respect to some non-empty subset of the image of $\phi$.
	\end{definition}

	\begin{proposition}\label{proposition:HellyTechnical}
		Let $T\ne K_1$ be a tree, $G$ a graph, $\{V_{\hat{x}}:\hat{x}\in V(T)\}$ a partition of $V(G)$ indexed by $V(T)$, and $k,\pow \ge 2$ a pair of integers.  If $G$ is $\c{F}_{T,k}^\pow$-free and if $\phi:V(T)\to V(G)$ is a canonical monomorphism, then every $T'\in \c{T}_\phi$ is leaf-cuttable and $\c{T}_\phi$ has the $k$-Helly property.
	\end{proposition}

	\begin{proof}[Proof of \Cref{theorem:treeBoundEG} assuming \Cref{proposition:HellyTechnical}]
		Let $G$ be an $n$-vertex $\c{F}_{T,k}^\pow$-free graph. By \Cref{lemma:canonicalReduction} there exists a partition of $V(G)$ into sets $\{V_{\hat{x}}:\hat{x}\in V(T)\}$ and a subset $\Phi\subseteq \mathrm{Mon}(T,G)$ with $|\Phi|=\Theta(\mathrm{mon}(T,G))$ such that $\phi(\hat{x})\in V_{\hat{x}}$ for all $\phi\in \Phi$ and $\hat{x}\in V(T)$, i.e.\ such that each $\phi\in \Phi$ is canonical.  Our ultimate goal will be to show that $\mathrm{mon}(T,G)=O(e(G)^{k-1})$, which by \Cref{observation monomorphisms same as copies} implies that $G$ contains at most this many copies of $T$ as desired. Since $|\Phi|=\Theta(\mathrm{mon}(T,G))$, it will suffice to prove $|\Phi|=O(e(G)^{k-1})$.  
		
		By Proposition~\ref{proposition:HellyTechnical} and \Cref{thm:Helly}, for every $\phi\in \Phi$, there exists a set of at most $k-1$ edges $\widehat{E}_\phi\sub E(T)$ such that every $T'\in \c{T}_\phi$ contains an edge of $\widehat{E}_\phi$, and without loss of generality we may assume that $\widehat{E}_\phi$ is non-empty (since if $\widehat{E}_\phi$ were empty we could add an arbitrary edge of $T\ne K_1$ to it while keeping its size $1\le k-1$).  Define 
		\[
		E_\phi:=\{\phi(\hat{x})\phi(\hat{y}):\hat{x}\hat{y}\in \widehat{E}_\phi\},
		\] 
		noting that $E_\phi$ is a non-empty set of at most $k-1$ edges of $E(G)$ since $\phi$ is a homomorphism, and that for every $T'\in \c{T}_\phi$ there exists an edge $uv\in E_\phi$ with $\hat{u}\hat{v}\in E(T')$. % and that $\chi(E_\phi)=E_\phi'$.
		
		We aim to show that for each non-empty set of at most $k-1$ edges $E\sub E(G)$, the number of $\phi\in \Phi$ with $E_{\phi}=E$ is $O(1)$, from which it will follow that $|\Phi|=O(e(G)^{k-1})$ since there are at most $e(G)^{k-1}$ choices for $E$.
		
		From now on we fix some $E$ as above.	Let $V_E\sub V(G)$ denote the set of vertices which are contained in at least one edge of $E$, let $T'_1,\ldots,T'_c$ denote the connected components of $T-\widehat{V}_E$, let $S_j\sub V_E$ denote the set of vertices $u\in V_E$ such that $\hat{u}$ is adjacent to a vertex of $T'_j$, and let $T''_j$ denote the subtree consisting of $T'_j$ and $\widehat{S}_j$.   We begin by observing some basic facts.
		
		\begin{claim}\label{cl:ETFacts}
			The following holds for each $j\ge 1$:
			\begin{enumerate}[label=(\roman*)]
				\item The set $S_j$ is non-empty,\label{claim ETFacts (i) non-empty S_j}
				\item Each $\hat{x}\in \widehat{S}_j$ is a leaf in $T''_j$\label{claim ETFacts (ii) each element of hat S_j is a leaf}
				\item No edge $uv\in E$ has $\hat{u}\hat{v}\in E(T''_j)$,\label{claim ETFacts (iii) no edge is in E(T''_j)}
				\item If $|\Psi(T''_j;S_j)|\ge \Pow$, then $T''_j$ is a highly-extendable subtree with respect to $S_j$.\label{claim ETFacts (iv) highly extentable}
			\end{enumerate}
		\end{claim}
		\begin{proof}
			For~\ref{claim ETFacts (i) non-empty S_j}, since $E$ is non-empty we have that $V_E$ is non-empty.  This means that each connected component of $T-\widehat{V}_E$ is a proper subgraph of $T$, and in particular the connected component $T'_j$ must have some vertex in $\widehat{V}_E$ which is adjacent to it in $T$, showing that $S_j$ is non-empty. 
			
			
			For~\ref{claim ETFacts (ii) each element of hat S_j is a leaf}, we have by definition that each $\hat{x}\in \widehat{S}_j$ is adjacent to at least one vertex $\hat{y}$ of $T'_j\sub T''_j$.  If $\hat{x}$ were also adjacent to some $\hat{z}\ne \hat{y}$ in $T'_j$, then since $T'_j$ is connected there would exist a cycle in $T''_j\sub T$ formed by taking the path from $\hat{y}$ to $\hat{z}$ in $T'_j$ together with the edges from $\hat{x}$ to $\hat{y},\hat{z}$, a contradiction to $T$ being a tree.  Similarly if there existed some $\hat{x}'\in \widehat{S}_j$ which $\hat{x}$ was adjacent to in $T''_j$, then since $\hat{x}'\in \widehat{S}_j$ there exists some neighbor $\hat{y}'$ of $\hat{x}'$ in $T'_j$, and now the (possibly empty) path from $\hat{y}'$ to $\hat{y}$ in $T'_j$ together with the vertices $\hat{x},\hat{x}'$ would create a cycle in $T'_j$, giving a contradiction.
			
			
			
			For~\ref{claim ETFacts (iii) no edge is in E(T''_j)}, if $uv\in E$, then $\hat{u},\hat{v}\not\in V(T'_j)\sub V(T)\sm \widehat{V}_E$ by definition.  As such, having $\hat{u}\hat{v}\in E(T_j'')$ would imply that $\hat{u}\in \widehat{S}_j$ is not a leaf (since $\hat{u}$ is adjacent to some vertex in $T'_j$ by definition of $S_j$ in addition to the vertex $\hat{v}$), a contradiction to what we proved in~\ref{claim ETFacts (ii) each element of hat S_j is a leaf}.
			
			For~\ref{claim ETFacts (iv) highly extentable}, condition~\ref{condition highly extendable definition (c) many canonical morphisms} of being highly-extendable holds by hypothesis, and condition~\ref{condition highly extendable definition (a) leaf of T'} holds by~\ref{claim ETFacts (ii) each element of hat S_j is a leaf}. For~\ref{condition highly extendable defintion (b) edges incident}, let $\hat{u}\hat{v}\in E(T)\setminus E(T_j'')$ be an edge incident to $T_j''$, say with $\hat{u}\in V(T_j'')$, $\hat{v}\not\in V(T_j'')$. Assume to the contrary that $\hat{u}\not\in \widehat{S}$, so $\hat{u}\in V(T_j')$. If $\hat{v}\in \widehat{V}_E$, then $\hat{v}\in \widehat{S}\subseteq V(T_j'')$ by definition, a contradiction. If on the other hand $\hat{v}\not\in \widehat{V}_E$, then actually $\hat{u}$ and $\hat{v}$ must be in the same component of $T-\widehat{V}_E$, namely $T'_j$, a contradiction to $\hat{u}\hat{v}\notin E(T''_j)$.
		\end{proof}
		We next observe that the situation in \Cref{cl:ETFacts}\ref{claim ETFacts (iv) highly extentable} never occurs in the cases we care about.
		
		
		
		\begin{claim}\label{cl:ETSmall}
			If there exists some $\phi\in \Phi$ such that $E=E_\phi$, then $|\Psi(T''_j;S_j)|<\Pow(T''_j,q)$ for all $j$.
		\end{claim}
		
		\begin{proof}
			If this were not the case for some $j$, then $T''_j$ would be a highly-extendable subtree with respect to $S_j$ by \Cref{cl:ETFacts}\ref{claim ETFacts (iv) highly extentable}.  Observe that having $E=E_\phi$ in particular implies (by definition of $E_\phi$) that $V_E$ is in the image of $\phi$. Since $S_j\sub V_E$ is a non-empty subset of the image of $\phi$ by \Cref{cl:ETFacts}\ref{claim ETFacts (i) non-empty S_j}, we conclude that $T''_j\in\c{T}_\phi$ by definition.  However, by definition of $E_\phi$ there must exist some edge $uv\in E_\phi$ with $\hat{u}\hat{v}\in E(T''_j)$, but this does not hold by \Cref{cl:ETFacts}\ref{claim ETFacts (iii) no edge is in E(T''_j)}.
			
		\end{proof}
		
		Let $\displaystyle\Pow=\max_{T'\sub T} \Pow(T',q)$.  We claim that the number of $\phi\in \Phi$ with $E=E_\phi$ is at most $\Pow^{v(T)}=O(1)$.  This is certainly true if no such $\phi$ exists, so we may assume that at least one such $\phi$ exists.  For such $\phi$, let $\phi_j:=\phi|_{V(T''_j)}$ for $j\ge 1$ and let $\phi_0:=\phi|_{\widehat{V}_E}$.  Observe that $\phi$ is uniquely determined by the vector $(\phi_0,\phi_1,\ldots,\phi_c)$ since every $\hat{x}\in V(T)$ is in either $\widehat{V}_E$ or one of the sets $V(T'_j)\subseteq V(T''_j)$ by definition of $T'_j$ being the connected components of $T-\widehat{V}_E$.  Also observe that $\phi_0$ is uniquely determined by $E$ since $\phi$ is canonical, and that each $\phi_j$ with $j\ge 1$ is an element of $\Psi(T''_j;S_j)$.  Thus by \Cref{cl:ETSmall}, there are at most $\Pow^{c}\le \Pow^{v(T)}$ choices of vectors $(\phi_0,\ldots,\phi_c)$, and hence at most this many choices of $\phi$, proving the claim.  It follows that $|\Phi|=O(e(G)^{k-1})$, giving the result.
	\end{proof}

\end{document}
